\chapter{Introducción}
La informática gráfica es un campo de las ciencias de la computación que se encarga de la creación, almacenamiento y visualización de imágenes en el ordenador. Este área utiliza algoritmos complejos para transformar los datos y modelos abstractos en los pixeles que acaban formando las imágenes en nuestras pantallas. 
Para esto hace uso de distintos campos, desde el álgebra y el cálculo para colocar los objetos en el espacio, hasta la óptica física para entender cómo la luz interactúa con las superficies y cómo interpreta nuestro cerebro los colores y el movimiento.\\

Principalmente el campo se divide en dos grandes bloques. Por un lado, está el modelado 3D que se encarga de definir la geometría y la estructura de los objetos, mientras que por el otro, se encuentra el proceso de renderizado (o visualización). Este último es encargado de calcular cómo afecta la luz a los modelos generados en el bloque anterior para crear la imagen final. Hoy en día, gracias a la enorme potencia de los equipos actuales, también se generan animaciones con físicas reales, visualizaciones complejas como las de datos médicos y gráficos impulsados directamente por la inteligencia artificial.\\

La informática gráfica tiene una gran variedad de aplicaciones que abarcan prácticamente todos los sectores de la industria e investigación. En arquitectura y diseño industrial es muy útil para proyectar estructuras y evaluar materiales y optimizar procesos, mientras que en medicina permite la exploración no invasiva a través de tomografías computarizadas visualizando el interior del cuerpo en tres dimensiones. También es fundamental para la simulación de fenómenos naturales complejos, analizar datos topográficos y estudiar la dinámica de fluidos.\\
 
En este último ámbito, la integración de Sistemas de Información Geográfica (SIG) permite procesar datos reales para generar una topografía avanzada y realista. Dichos datos se consiguen principalmente de nubes de puntos LiDAR (Light Detection and Ranging), que son conjuntos masivos de coordenadas que se han obtenido al mapear la superficie terrestre con alta precisión. A partir del filtrado de estos puntos se pueden obtener Modelos Digitales de Elevación (MDE) que son matrices bidimensionales donde cada píxel almacena el valor exacto de la altitud del terreno, pudiendo crear con ellos mapas de alturas (heightmaps). Estas representaciones son vitales para el modelado de simulaciones complejas de flujos de fluidos, meteorología o la propagación de desastres naturales ya que proporcionan el relieve tridimensional exacto sobre el que interactúan los modelos físicos.\\

Para llevar a cabo estas representaciones, tradicionalmente se ha optado por el modelado mediante mallas, que consisten en una estructura  formada por vértices colocados en el espacio y que se unen mediante aristas creando polígonos (generalmente triángulos) que acaban dando forma a un modelo el cual está vacío en su interior. Posteriormente la renderización de estas mallas se realiza mediante la rasterización, que es el proceso algorítmico que proyecta estos triángulos tridimensionales sobre un plano bidimensional, pudiendose representar así en los píxeles de una pantalla.\\

Este enfoque es muy eficiente para representar elementos sólidos con limites bien definidos, como el relieve de una montaña o un edificio pero cuando se trata de modelar fenómenos amorfos y caóticos como el fuego las mallas poligonales no son funcionales ya que estos elementos no tienen una superficie sólida y estática. Para solucionar esto, se decide usar el modelado volumétrico mediante vóxeles, que utiliza una estructura tridimensional de celdas para representar no solo la frontera, sino el interior de un modelo, donde cada vóxel almacena propiedades físicas continuas, como la densidad, la temperatura o la concentración de combustible.\\

Al igual que cambia el modelado el algoritmo de renderizado también cambia cuando nos referimos a modelos volumétricos como el del fuego y pasamos de usar la rasterización a implementar técnicas de Renderizado Volumétrico Directo (DVR).  A diferencia de los métodos indirectos que extraen mallas poligonales intermedias para representar la capa exterior de un volumen, el DVR permite procesar y visualizar toda la información de forma directa sin depender de geometrías exactas.\\

Existen varias técnicas de DVR que, aunque fueron pioneras, presentan limitaciones para simulaciones dinámicas: 
\begin{itemize}
    \item \textbf{Splatting (Proyección de huellas):} En esta técnica, el algoritmo procesa el volumen proyectando cada vóxel individualmente sobre el plano de la imagen. A medida que estas huellas se van acumulando en la pantalla, sus colores y transparencias se suman para formar la imagen final. Aunque es un método rápido, presenta problemas de solapamiento y tiende a generar imágenes borrosas cuando se mueve la cámara.
    
    \item \textbf{Renderizado basado en texturas (Texture-based Volume Rendering):} Este método divide el volumen en planos que se alinean de forma paralela a la cámara y se les asigna la información de los vóxeles como si fuera una textura 2D. Tienen dos problemas principales, por una parte si se rota la cámara todo el sistema de planos debe recalcularse para volver a alinearse con esta, y por otro lado un número bajo de láminas produce un efecto visual escalonado mientras que un número alto dispara el consumo de memoria.
\end{itemize}

Debido a estos problemas, las dos técnicas podrían ser consideradas para la representación de volúmenes estáticos, pero para simulaciones dinámicas como las de un incendio forestal no serían lo suficientemente flexibles.\\ 

Para el desarrollo de estos y poder aprovechar al máximo la capacidad de procesamiento paralelo de las tarjetas gráficas modernas, la técnica elegida para este simulador es el trazado de rayos volumétrico(\textit{Raymarching}). Este algoritmo lanza rayos virtuales desde la cámara hacia la cuadrícula 3D, avanzando a través del volumen en pequeños pasos, donde en cada uno de ellos evalúa los datos del vóxel y calcula cómo la luz se acumula, se bloquea o se dispersa al atravesarlo.\\

Además de la visualización de los volúmenes, en un simulador hay que añadirle el coste computacional del cálculo de las físicas continuas. Para ello normalmente se requiere resolver complejas ecuaciones para cada voxel que se actualice en sistemas de autómatas celulares. Realizar estos cálculos  en la unidad central de procesamiento (CPU) resultaría prácticamente imposible y rompería el simulador.\\



